\documentclass{article}


\usepackage{PRIMEarxiv}

\usepackage[utf8]{inputenc} % allow utf-8 input
\usepackage[T1]{fontenc}    % use 8-bit T1 fonts
\usepackage{hyperref}       % hyperlinks
\usepackage{url}            % simple URL typesetting
\usepackage{booktabs}       % professional-quality tables
\usepackage{amsfonts}       % blackboard math symbols
\usepackage{nicefrac}       % compact symbols for 1/2, etc.
\usepackage{microtype}      % microtypography
\usepackage{fancyhdr}       % header
\usepackage{graphicx}       % graphics
\usepackage{amsmath}
\graphicspath{{media/}}     % organize your images and other figures under media/ folder

%Header
\pagestyle{fancy}
\thispagestyle{empty}
\rhead{ \textit{ }} 

% Update your Headers here
\fancyhead[LO]{Identifying Landing Sites for Mars Landers}
% \fancyhead[RE]{Firstauthor and Secondauthor} % Firstauthor et al. if more than 2 - must use \documentclass[twoside]{article}



  
%% Title
\title{Characterizing Single-Frame Safe-Zone Estimation for Mars Landers}

\author{
  Kalp Shah \\
  Department of Electrical Engineering \\
  Indian Institute of Technology Gandhinagar \\
  \texttt{kalp.shah@iitgn.ac.in} \\
  %% examples of more authors
   \And
  Pramod Mishra \\
  Department of Electrical Engineering \\
  Indian Institute of Technology Gandhinagar \\
  \texttt{pramod.25330013@iitgn.ac.in} \\
  \And
  Shanmuganathan Raman \\
  Department of Computer Science and Engineering \\
  Indian Institute of Technology Gandhinagar \\
  \texttt{shanmuga@iitgn.ac.in} \\
  %% \AND
  %% Coauthor \\
  %% Affiliation \\
  %% Address \\
  %% \texttt{email} \\
  %% \And
  %% Coauthor \\
  %% Affiliation \\
  %% Address \\
  %% \texttt{email} \\
  %% \And
  %% Coauthor \\
  %% Affiliation \\
  %% Address \\
  %% \texttt{email} \\
}


\begin{document}
\maketitle


\begin{abstract}
Safe autonomous landing on planetary surfaces requires quick, accurate decisions under tight computational limits. This paper characterizes a deterministic computer vision pipeline for identifying safe landing zones from high-resolution single-frame nadir orbital-style imagery. The pipeline fuses gradient-based intensity edge evaluation, contour-based shadow masking, and FFT-based image texture classification into a unified hazard score. A maximal-rectangle algorithm then extracts the largest contiguous safe landing zone. Evaluated on a synthetic dataset utilizing DTM-derived proxy labels (safe = slope $\le$ 15 degrees), the system achieves a proxy-label F1 score of 0.025 [0.016, 0.035], which is marginally above a random baseline (0.020). The low performance is primarily due to the texture module being overly sensitive to high-frequency noise, combined with the difficulty of predicting true geometric slope purely from single-frame intensity without a dedicated shape-from-shading or photometric stereo model. The pipeline's measured execution time is 343.5 ms on a 1024x1024 image on standard host hardware. The study demonstrates the limitations of purely intensity-based thresholding for geometric hazard avoidance and underscores the need for geometric proxies.
\end{abstract}


% keywords can be removed
\keywords{Computer Vision \and Mars Landers \and Autonomous Landing \and Hazard Detection \and Fast Fourier Transform}


\section{Introduction}
Autonomous landing on Mars requires a spacecraft to identify suitable terrain without continuous control from Earth. During terminal descent, communication delays make ground-in-the-loop decisions impractical, requiring onboard systems to assess terrain from sensor data.

Landing suitability is strictly bounded by mission-specific engineering constraints \cite{golombek2012selection}, relying on terrain properties including slope, surface roughness, shadows, and local surface texture. Existing approaches use range sensors \cite{amzajerdian2013alhat} or learning-based vision methods \cite{moghe2020deep}. Deep learning methods can provide strong visual representations, but their computational and power requirements can be challenging to accommodate on some radiation-hardened processors (e.g., the RAD750 \cite{berger2001rad750}) used for spaceflight, though the success of Mars 2020 Terrain Relative Navigation and the FPGA-based Lander Vision System demonstrate ongoing progress in flight computing \cite{johnson2020mars2020}. However, these processors still have substantially tighter power, memory, and compute constraints than modern terrestrial hardware, motivating lightweight deterministic approaches.

The main difficulty is the lack of ground-truth landing-safety annotations. Unlike conventional vision tasks, there is no established pixel-level label indicating whether a region of Martian terrain is physically safe for landing. This work therefore treats landing safety as an estimation problem: observable visual features are combined to produce an approximate map of terrain suitability.

We propose a deterministic computer-vision pipeline that analyzes image gradients, shadowed regions, and spatial-frequency characteristics of the terrain. These signals are fused into a safety map, from which a candidate landing region is selected using a maximal-rectangle algorithm. The resulting map is a heuristic estimate of landing suitability, not a calibrated probability or ground-truth classification.

\section{Related Work}
Autonomous planetary landing relies on robust hazard detection. Traditional systems, such as the Autonomous Landing Hazard Avoidance Technology (ALHAT), heavily utilized LiDAR for generating dense 3D elevation maps \cite{amzajerdian2013alhat}. To bypass the size, weight, and power (SWaP) penalties of active sensors, classical vision-based hazard detection and shadow-based methods were developed to infer safety directly from passive descent imagery \cite{restrepo2020next}. More recently, deep learning approaches, particularly Convolutional Neural Networks (CNNs), have been applied to identify craters and classify hazardous lunar terrain \cite{moghe2020deep}. However, due to the tight resource constraints of typical spaceflight hardware, there remains an interest in lightweight, classical computer vision pipelines that do not require extensive offline training or specialized AI accelerators.

\section{Methodology}
\label{sec:headings}
The proposed pipeline processes input terrain imagery through four stages to determine the optimal landing zone. It uses various techniques to independently characterize intensity gradients, shadows, and fine-scale roughness (texture), and then blends the information obtained from
individual techniques to generate a safety map pixel-wise. The
selection of the threshold is empirical and is derived from a
Martian image dataset available.

\subsection{Gradient Hazard Scoring}
The first module quantifies intensity gradient and roughness at the pixel level by computing the spatial gradient of the input grayscale image $I(x,y)$. Horizontal and vertical partial derivatives are obtained via the Sobel operator\cite{sobel1968} with a standard $3 \times 3$ kernel size:
\begin{equation}
    I_x = \frac{\partial I}{\partial x}, \qquad I_y = \frac{\partial I}{\partial y}.
\end{equation}
The gradient magnitude at each pixel is then defined as
\begin{equation}
    g(x,y) = \sqrt{I_x^2 + I_y^2},
\end{equation}
where large values correspond to abrupt intensity transitions associated with steep slopes, rock edges, or crater rims, which are features that are hazardous for landing. To produce a scene-adaptive, bounded hazard map, $g(x,y)$ is normalized by its 99th percentile $g_{p99}$ computed over all non-zero pixels:
\begin{equation}
    G(x,y) = \min\!\left(1,\; \frac{g(x,y)}{g_{p99}}\right).
\end{equation}
Using the 99th percentile rather than the global maximum suppresses the influence of outlier high-gradient pixels, ensuring robustness to rare extreme values. The resulting map $G(x,y) \in [0,1]$ is the primary hazard score, with values near 1.0 indicating highly hazardous terrain and near 0.0 indicating smooth, potentially safe terrain.

\subsection{Shadow Augmentation}
The gradient map alone does not capture all hazardous regions. The broad, flat interiors of craters and shadow-filled depressions exhibit low internal gradients, and are therefore incorrectly assigned low hazard scores, despite representing dangerous terrain. The second module corrects this systematic blind spot through morphological contour analysis.
 
A binary inverse threshold is applied to the input grayscale image at intensity level $\tau = 40$, extracting dark connected regions that correspond to shadows and crater interiors. External contours are then identified; contours whose enclosed area falls below $A_{\min} = 500$ pixels are discarded to eliminate noise and spurious small dark patches. For every retained contour, its enclosed interior is flood-filled, and the hazard score for all pixels within the fill is forced to the maximum value:
\begin{equation}
    G(x,y) \leftarrow 1.0 \quad \forall\,(x,y) \in \text{filled contour region}.
\end{equation}
The physical rationale is that shadows adjacent to mountains, fault scarps, and crater walls are reliable proxies for steep, unsafe terrain. By overriding the gradient score inside large dark contours, the module compensates for the low-gradient interior while preserving the spatial boundary information inherent in the contour geometry. The two key parameters, the intensity threshold $\tau = 40$ and the minimum contour area $A_{\min} = 500 pixels$, were chosen empirically to balance sensitivity and specificity under typical high-resolution nadir orbital-style frames at 1 m/px GSD illumination conditions.

\subsection{FFT-based Texture-Energy Score}
Surface roughness is an important terrain property that cannot be inferred from slope alone. Small, distributed hazards such as rock fields or patterned ground may lack steep macro-gradients while still presenting significant risks to lander stability. Shepard et al. \cite{shepard2001roughness} demonstrate that natural planetary terrain exhibits scale-dependent variations in surface topography. While they note that spatial methods often yield cleaner roughness parameters than power spectral methods, a localized Fourier-based analysis offers an efficient computational proxy to quantify the distribution of spatial-frequency energy as an indicator of fine-scale terrain roughness. A Fourier-based analysis is therefore used to quantify the distribution of spatial-frequency energy as an indicator of fine-scale terrain roughness.
 
The augmented image is partitioned into non-overlapping $32 \times 32$ pixel blocks. For each block $B_k$, its mean intensity is first subtracted to remove the DC component, and a centered two-dimensional discrete Fourier transform (DFT) is computed to obtain the spectral magnitude $|\hat{B}_k(u,v)|$. The spectral plane is then partitioned into a low-frequency region (radial frequency $\rho \leq r_0$) and a high-frequency region ($\rho > r_0$), where $\rho = \sqrt{u^2 + v^2}$ and $r_0 = 6$ pixels in the frequency domain. The corresponding energy terms are
\begin{equation}
    E_{\text{low}}  = \sum_{\rho \leq r_0} |\hat{B}_k(u,v)|^2, \qquad
    E_{\text{high}} = \sum_{\rho >   r_0} |\hat{B}_k(u,v)|^2,
\end{equation}
and the texture-energy ratio is defined as
\begin{equation}
    R = \frac{E_{\text{high}}}{E_{\text{low}} + \varepsilon},
\end{equation}
where $\varepsilon = 10^{-8}$ prevents division by zero. A block is classified as \emph{Rough} if $R > 1.0$ and as \emph{Smooth} otherwise. Rough terrain features, such as rock fields and small craters, produce rich high-frequency texture signatures in nadir orbital-style frames at 1 m/px GSD. Conversely, smooth, safe terrain is dominated by low-frequency variations corresponding to broad brightness gradients. It must be noted that while this module serves as a proxy for topographic roughness, it fundamentally measures image texture, which can also be influenced by surface albedo variations, illumination changes, and sensor noise. Each block is mapped to a terrain safety score that has been chosen empirically after visual inspection:
\begin{equation}
    \text{terrainSafetyScore} =
    \begin{cases}
        0.1 & \text{if Rough} \; (R > 1.0), \\
        1.0 & \text{if Smooth} \; (R \leq 1.0).
    \end{cases}
\end{equation}
The block-level classification is up-sampled to pixel resolution by assigning every pixel the score of its parent block, yielding a spatially registered texture-energy map aligned with the gradient hazard map from the previous modules. Because $0.1$ is strictly less than the fusion threshold $\theta = 0.35$, classifying a block as Rough acts as a hard veto, unconditionally excluding highly textured regions from the final safe mask.

\subsection{Safe-Region Fusion and Landing Box Selection}
The final module integrates the gradient hazard map and the texture-energy score into a single safe-landing score, then identifies the largest viable landing rectangle within the fused safe region.
 
For each pixel, the composite safe score is computed as
\begin{equation}
    s(x,y) = \bigl(1 - G(x,y)\bigr) \times \text{terrainSafetyScore}(x,y).
\end{equation}
The factor $(1-G)$ converts the hazard score into a safety score: pixels with high gradient hazard contribute minimally regardless of their texture classification, while smooth, low-texture pixels receive scores close to 1.0. A pixel is labeled safe if $s(x,y) > \theta$, where the threshold $\theta = 0.35$ was selected to admit the majority of genuinely safe terrain while rejecting borderline regions. The binary safe mask is processed with connected-component analysis to remove isolated small patches below a minimum area threshold, retaining only spatially contiguous landing-capable zones.
 
The recommended landing zone is then determined by applying the maximal-rectangle algorithm to the safe mask. This classic histogram-based algorithm \cite{vandevoorde1996maximal}, operating in $O(W)$ time per row and $O(HW)$ time overall, where $W$ and $H$ are the image width and height, identifies the largest axis-aligned rectangle that lies entirely within the safe region. The output is a bounding box parameterized by its top-left corner $(x_0, y_0)$ and dimensions $(w, h)$ in image coordinates. This rectangle is the recommended lander deployment zone, balancing landing-area maximization with safety considerations derived from the multimodal terrain analysis.


\begin{table}[htbp]
\centering
\caption{Pipeline Parameters and Selection Methodology}
\begin{tabular}{@{}lll@{}}
\toprule
\textbf{Parameter} & \textbf{Value} & \textbf{Selection Method} \\ \midrule
Sobel Kernel & $3 \times 3$ & Standard practice \cite{sobel1968} \\
Gradient Normalization ($g_{p99}$) & 99th percentile & Empirical (outlier suppression) \\
Shadow Intensity Threshold ($\tau$) & 40 (8-bit) & Empirical visual inspection \\
Minimum Contour Area ($A_{\min}$) & 500 px & Empirical noise filtering \\
FFT Block Size & $32 \times 32$ px & Localization/resolution trade-off \\
Frequency Threshold ($r_0$) & 6 px & Empirical \\
Texture Regularizer ($\varepsilon$) & $10^{-8}$ & Numerical stability \\
Texture-Energy Ratio Threshold & 1.0 & Empirical separation of texture \\
Safe Fusion Cutoff ($\theta$) & 0.35 & Empirical safe zone cutoff \\\\\\
Minimum Component Area & 1500 px & Empirical (noise filtering) \\\\
\bottomrule
\end{tabular}
\label{tab:parameters}
\end{table}

\section{Results}
\label{sec:results}
The pipeline was evaluated on a synthetically generated dataset of 80 TEST tiles mapping 1 m/px nadir imagery to DTMs. Proxy labels were constructed using a 5m baseline where slope $> 15^\circ$ is proxy-unsafe. 

\subsection{Quantitative Metrics}
The full pipeline (v0) achieves a pixel-level macro-averaged F1 score of 0.025 [0.016, 0.035] for the safe class, with an unsafe recall of 0.935. In comparison, a random baseline (B5, matched safe fraction) achieves F1=0.020, and a local standard deviation baseline (B1) achieves F1=0.720. Gradient-only ablation (A1) yields F1=0.897. The results demonstrate that while the FFT and gradient modules correlate with rough terrain, strict static thresholding makes the pipeline overly pessimistic, rejecting the vast majority of the terrain. The Spearman correlation between the gradient hazard $G$ and true DTM slope is 0.340, confirming a moderate signal that is lost during binary thresholding. 

\subsection{Runtime Analysis}
The total pipeline execution time on a single thread of host hardware (Intel/AMD) for a 1024x1024 image is 343.5 ms (median over 50 runs after warm-up). The most expensive component is the FFT block processing.


\section{Limitations}
This evaluation highlights several limitations. First, intensity gradients correlate only moderately with physical slope due to illumination dependence. Second, the use of proxy labels means physical safety is approximated. Third, the pipeline assumes nadir orbital-style imagery, lacking tests on oblique descent imagery. Finally, the system lacks a proximity penalty, and execution times are measured only on host hardware, not flight processors.

\section{Conclusion}
This paper characterized a classical, deterministic computer vision pipeline for single-frame safe-zone estimation from nadir orbital-style imagery. By fusing Sobel-based intensity gradient scoring, contour-based shadow augmentation, and an FFT-based amplitude-invariant texture-energy score, the pipeline produces a per-pixel safety map, from which the largest viable landing rectangle is extracted via a maximal-rectangle algorithm \cite{vandevoorde1996maximal}. Quantitative evaluation on a synthetic DTM proxy dataset revealed that the system achieves a proxy-label F1 score of 0.025. While the pipeline executes in 343.5 ms on host hardware, the results indicate that purely intensity-based classical thresholding struggles to reliably isolate physically safe geometric zones without explicit distance penalties or photometric priors.

\section*{Acknowledgments}

The authors acknowledge the Indian Institute of Technology, Gandhinagar, for providing the infrastructure and resources that supported this work.

%Bibliography
\bibliographystyle{unsrt}  
\bibliography{references}  


\end{document}
